\section{Integer-shift Stieltjes correlations from a regularized depth-two kernel}
\label{sec:stieltjes-generator}

The tail generator admits a useful singular-order limit.  Write the
Laurent expansion of the Hurwitz zeta function at one as
\begin{equation}
 \zeta(1+\alpha,a)=\frac1\alpha+
 \sum_{r=0}^{\infty}\frac{(-1)^r}{r!}\gamma_r(a)\alpha^r,
 \label{eq:stieltjes-laurent}
\end{equation}
where $\gamma_r(1)=\gamma_r$ and $\gamma_0=\gamma$; see, for example,
\cite{DLMF,Coffey}.  At integer shifts,
\begin{equation}
 \gamma_r(n+1)=\gamma_r-
 \sum_{m=1}^n\frac{\log^r m}{m}.
 \label{eq:stieltjes-shift}
\end{equation}
This follows by comparing \eqref{eq:stieltjes-laurent} in
$\zeta(1+\alpha,n+1)=\zeta(1+\alpha)-\sum_{m=1}^nm^{-1-\alpha}$.

\subsection{Single-index order derivatives}

\begin{proposition}[Integer-shift Stieltjes generator]
\label{prop:single-stieltjes-generator}
For every $r\ge0$ and $|z|<1$,
\begin{equation}
 \boxed{\displaystyle
 \sum_{n=0}^{\infty}\gamma_r(n+1)z^n
 =\frac{\gamma_r-(-1)^r
 \left.\partial_s^r\Li_s(z)\right|_{s=1}}{1-z}.}
 \label{eq:single-stieltjes-generator}
\end{equation}
\end{proposition}

\begin{proof}
Sum \eqref{eq:stieltjes-shift} over $n$ and interchange the absolutely
convergent sums.  The logarithmic Dirichlet series is
\[
 \sum_{m\ge1}\frac{z^m\log^r m}{m}
 =(-1)^r\left.\partial_s^r\Li_s(z)\right|_{s=1}.
\]
\end{proof}

The point of the next theorem is that products do not require an
uncontrolled double convolution of the shifted constants.  A single
pole-subtracted depth-two kernel generates all of them.

\subsection{The bivariate regularized kernel}

For $|z|<1$, define
\begin{align}
 \K(\alpha,\beta;z)={}&
 \Li_{1+\alpha,1+\beta}(1,z)
       -\frac{\Li_{1+\beta}(z)}{\alpha}
 \notag\\
 &+\Li_{1+\beta,1+\alpha}(1,z)
       -\frac{\Li_{1+\alpha}(z)}{\beta}
 +\Li_{2+\alpha+\beta}(z).
 \label{eq:K-kernel}
\end{align}
Although the displayed terms have poles, their specified differences
are regular.

\begin{lemma}[Holomorphy of the pole subtraction]
\label{lem:K-holomorphic}
For every fixed $|z|<1$, $\K(\alpha,\beta;z)$ extends holomorphically
to a neighborhood of $(0,0)$.
\end{lemma}

\begin{proof}
For $\Re\alpha>0$,
\begin{align*}
 \Li_{1+\alpha,1+\beta}(1,z)
 &=\sum_{n\ge1}\frac{z^n}{n^{1+\beta}}
   \zeta(1+\alpha,n+1),\\
 \Li_{1+\alpha,1+\beta}(1,z)
 -\frac{\Li_{1+\beta}(z)}{\alpha}
 &=\sum_{n\ge1}\frac{z^n}{n^{1+\beta}}
 \left(\zeta(1+\alpha,n+1)-\frac1\alpha\right).
\end{align*}
The parenthesis is holomorphic at $\alpha=0$.  On compact subsets of a
small bidisk it grows at most polynomially in $\log(n+1)$, whereas
$|z|^n$ decays exponentially.  The series is therefore locally
normally convergent and holomorphic in both variables.  The second
pole-subtracted term is symmetric, and the ordinary polylogarithm in
\eqref{eq:K-kernel} is already holomorphic in its order near two.
\end{proof}

Put
\begin{equation}
 G(\alpha)=\zeta(1+\alpha)-\frac1\alpha
 =\sum_{r\ge0}\frac{(-1)^r\gamma_r}{r!}\alpha^r.
 \label{eq:G-alpha}
\end{equation}

\begin{theorem}[All-index Stieltjes correlation generator]
\label{thm:stieltjes-correlation-generator}
For $|z|<1$ and $(\alpha,\beta)$ near $(0,0)$,
\begin{equation}
 \boxed{\displaystyle
 \sum_{n=0}^{\infty}
 \left(\zeta(1+\alpha,n+1)-\frac1\alpha\right)
 \left(\zeta(1+\beta,n+1)-\frac1\beta\right)z^n
 =\frac{G(\alpha)G(\beta)-\K(\alpha,\beta;z)}{1-z}.}
 \label{eq:parameter-correlation}
\end{equation}
Consequently, for all $r,s\ge0$,
\begin{equation}
 \boxed{\displaystyle
 \sum_{n=0}^{\infty}\gamma_r(n+1)\gamma_s(n+1)z^n
 =\frac{\gamma_r\gamma_s
 -(-1)^{r+s}r!s!\,[\alpha^r\beta^s]
       \K(\alpha,\beta;z)}{1-z}.}
 \label{eq:stieltjes-correlation}
\end{equation}
\end{theorem}

\begin{proof}
Apply \cref{thm:two-tail-generator} formally first in the half-plane
$\Re\alpha,\Re\beta>0$, with $p=1+\alpha$ and $q=1+\beta$.  Also use
the one-tail identity
\[
 \sum_{n\ge0}\zeta(1+\alpha,n+1)z^n
 =\frac{\zeta(1+\alpha)-\Li_{1+\alpha}(z)}{1-z}.
\]
Subtract $1/\alpha$ from the first factor and $1/\beta$ from the
second.  The zeta terms combine into $G(\alpha)G(\beta)$, while the
remaining polylogarithmic terms combine into $-\K$.  This proves
\eqref{eq:parameter-correlation} in a nonempty domain.  The holomorphic
extensions from \cref{lem:K-holomorphic} continue it to a neighborhood
of the origin.  Coefficient extraction using
\eqref{eq:stieltjes-laurent} gives
\eqref{eq:stieltjes-correlation}.
\end{proof}

\subsection{Digamma square and an anchored antiderivative}

Since $\gamma_0(n+1)=\gamma-H_n=-\psi(n+1)$, the lowest coefficient of
\cref{thm:stieltjes-correlation-generator} simplifies completely.

\begin{corollary}[Digamma-square generating function]
\label{cor:digamma-square}
For $|z|<1$,
\begin{equation}
 \boxed{\displaystyle
 \sum_{n=0}^{\infty}(\gamma-H_n)^2z^n
 =\sum_{n=0}^{\infty}\psi(n+1)^2z^n
 =\frac{(\gamma+\log(1-z))^2+\Li_2(z)}{1-z}.}
 \label{eq:digamma-square}
\end{equation}
\end{corollary}

\begin{proof}
The elementary generating functions
\[
 \sum_{n\ge0}H_nz^n=-\frac{\log(1-z)}{1-z},
 \qquad
 \sum_{n\ge0}H_n^2z^n
 =\frac{\Li_2(z)+\log^2(1-z)}{1-z}
\]
follow by summing the defining single and double harmonic sums.  Expand
$(\gamma-H_n)^2$ and combine them.
\end{proof}

Termwise integration provides the anchored primitive requested by the
factor $1/(n+1)$.

\begin{theorem}[Anchored digamma-square primitive]
\label{thm:digamma-square-primitive}
For $0\le z<1$,
\begin{align}
 \sum_{n=0}^{\infty}\frac{(\gamma-H_n)^2}{n+1}z^{n+1}
 ={}&\frac{\gamma^3-(\gamma+\log(1-z))^3}{3}
 -\Li_2(z)\log(1-z)-I(z),
 \label{eq:digamma-primitive}\\
 I(z):={}&\int_0^z\frac{\log^2(1-t)}t\dd t
 \notag\\
 ={}&2\zeta(3)+\log^2(1-z)\log z
 +2\log(1-z)\Li_2(1-z)-2\Li_3(1-z).
 \label{eq:I-closed}
\end{align}
In particular,
\begin{equation}
 \boxed{\displaystyle
 \sum_{n=0}^{\infty}
 \frac{(\gamma-H_n)^2}{(n+1)2^{n+1}}
 =\frac{\gamma^3-(\gamma-\log2)^3}{3}
 +\frac{\pi^2\log2}{12}-\frac{\log^32}{6}-\frac{\zeta(3)}4.}
 \label{eq:half-digamma-primitive}
\end{equation}
\end{theorem}

\begin{proof}
Integrate \eqref{eq:digamma-square} from $0$ to $z$.  The square term
is integrated by the substitution $u=\gamma+\log(1-t)$.  Integration
by parts gives
\[
 \int_0^z\frac{\Li_2(t)}{1-t}\dd t
 =-\Li_2(z)\log(1-z)-I(z).
\]
Differentiating the right side of \eqref{eq:I-closed} and checking its
zero value at $z=0$ proves the closed form.  At $z=1/2$, substitute the
classical evaluations
\[
 \Li_2(\tfrac12)=\frac{\pi^2}{12}-\frac{\log^22}{2},
 \qquad
 \Li_3(\tfrac12)=\frac78\zeta(3)-\frac{\pi^2\log2}{12}
                  +\frac{\log^32}{6}
\]
and simplify.
\end{proof}

More generally, integrating \eqref{eq:stieltjes-correlation} gives an
anchored antiderivative for every pair of indices:
\begin{equation}
 \sum_{n=0}^{\infty}
 \frac{\gamma_r(n+1)\gamma_s(n+1)}{n+1}z^{n+1}
 =\int_0^z
 \frac{\gamma_r\gamma_s-(-1)^{r+s}r!s!
 [\alpha^r\beta^s]\K(\alpha,\beta;t)}{1-t}\dd t.
 \label{eq:all-index-primitive}
\end{equation}
The point is not merely the integral notation: the integrand is a
finite coefficient of the explicitly specified depth-two kernel
\eqref{eq:K-kernel}, so order derivatives and anchored primitives are
controlled by one common analytic object.
